\section{Experiments}
\label{sec:experiments}

We evaluate the proposed pipeline on three classes of problems: a
two-dimensional incompressible Navier--Stokes benchmark (Section~\ref{sec:tgv}),
two canonical viscous-flow profiles (Section~\ref{sec:couette_poiseuille}), and
a pair of algebraic physical laws recovered from noisy data
(Section~\ref{sec:algebraic}).  Section~\ref{sec:ablation} presents ablation
studies isolating the contribution of each design choice.  All experiments use
PyTorch~2.0+ on a single CPU; the Taylor--Green search completes in under
four minutes.

% ===========================================================================
\subsection{Taylor--Green Vortex (2-D Incompressible Navier--Stokes)}
\label{sec:tgv}

The Taylor--Green vortex is an exact, time-decaying solution to the
two-dimensional incompressible Navier--Stokes equations:
\begin{align}
  u(x,y,t) &= \phantom{-}\sin(x)\cos(y)\,\exp(-2\nu t), \label{eq:tgv_u}\\
  v(x,y,t) &= -\cos(x)\sin(y)\,\exp(-2\nu t), \label{eq:tgv_v}\\
  p(x,y,t) &= -\tfrac{1}{4}\bigl[\cos(2x)+\cos(2y)\bigr]\exp(-4\nu t),
              \label{eq:tgv_p}
\end{align}
where $\nu=0.1$ is the kinematic viscosity.  Because the velocity field is a
product of separable, axis-aligned elementary functions, this problem provides
an ideal test of whether the \textsc{ComposeHead} can recover exact symbolic
structure from data alone.

\paragraph{Setup.}
We sample $n=3{,}000$ training points uniformly over
$x,y\in[0,2\pi]$, $t\in[0,1]$ and evaluate both velocity components at the
exact solution~\eqref{eq:tgv_u}--\eqref{eq:tgv_v}.  The function basis is
$\mathcal{F}=\{\sin,\cos,\exp\}$, yielding $|\mathcal{F}|^3=27$ triple
separable candidates of the form $f(x)\,g(y)\,h(t)$.  Each candidate is
trained independently for 2{,}000 steps with Adam (learning rate $0.01$,
gradient clipping at max-norm~$1.0$).  The best candidate for each velocity
component is then fine-tuned for an additional 5{,}000 steps at learning rate
$0.005$.

\paragraph{Search results.}
For the $u$-component, the search identifies
$\sin(\cdot)\!\cdot\!\sin(\cdot)\!\cdot\!\exp(\cdot)$ as the lowest-loss
candidate (loss $<10^{-8}$); for $v$, the winner is
$\cos(\cdot)\!\cdot\!\cos(\cdot)\!\cdot\!\exp(\cdot)$.  These selections appear
to differ from the ground truth, which contains $\sin\!\cdot\!\cos$ and
$\cos\!\cdot\!\sin$ factors, respectively.  The discrepancy is resolved by
observing that the optimizer absorbs phase shifts into the bias parameters:
$\sin(x - \pi)\sin(y - \pi/2)\exp(at+b) = \sin(x)\cos(y)\exp(at)\exp(b)$
up to sign, so the learned representation is equivalent to the true solution
modulo a reparameterization that the normalization step removes.

\paragraph{Raw learned parameters.}
After fine-tuning, the combined MSE loss reaches $\sim\!10^{-10}$.  The raw
parameters before normalization are as follows:
\begin{itemize}
  \item $u$: $\text{coeff}=0.540$;\;
        $\sin_0$: $a\!=\!1.000$, $b\!=\!{-3.142}\;(\approx -\pi)$;\;
        $\sin_1$: $a\!=\!1.000$, $b\!=\!{-1.571}\;(\approx -\pi/2)$;\;
        $\exp_2$: $a\!=\!{-0.200}$, $b\!=\!0.616$.
  \item $v$: $\text{coeff}=0.594$;\;
        $\cos_0$: $a\!=\!1.000$, $b\!=\!0.000$;\;
        $\cos_1$: $a\!=\!1.000$, $b\!=\!1.571\;(\approx \pi/2)$;\;
        $\exp_2$: $a\!=\!{-0.200}$, $b\!=\!0.522$.
\end{itemize}

\paragraph{Normalization and snapping.}
The normalization step (i)~absorbs exponential biases into the coefficient
($e^{0.616}\approx 1.851$, yielding $0.540 \times 1.851 \approx 1.000$),
(ii)~canonicalizes trigonometric phases by applying the identities
$\sin(x-\pi)=-\sin(x)$ and $\sin(y-\pi/2)=\cos(y-\pi)\mapsto -\cos(y)$
together with the sign-absorption rule for $\sin$, and
(iii)~reduces all bias terms to zero.
After normalization, coefficient snapping with tolerance~$0.05$ yields the
exact closed-form expressions:
\begin{equation}
  \hat{u} = \sin(x)\cos(y)\exp(-t/5), \qquad
  \hat{v} = -\cos(x)\sin(y)\exp(-t/5).
  \label{eq:recovered}
\end{equation}
The snap error---defined as the mean absolute difference between the snapped
model and the ground truth on 5{,}000 held-out points---is $7\times10^{-8}$,
consistent with machine precision.

\paragraph{PDE residual verification.}
We verify the recovered solution against the governing equations.  The
continuity residual $|\nabla\!\cdot\!\mathbf{u}|$ has mean value
$3\times10^{-4}$, attributable to finite-difference approximation of the
divergence.  The momentum residuals are non-zero because the pipeline does not
model the pressure field; the residual in each momentum equation equals
$\nabla p$, which is consistent with the pressure
gradient implied by~\eqref{eq:tgv_p}.

\paragraph{Comparison with baselines.}
Table~\ref{tab:tgv_comparison} compares four approaches on the same
Taylor--Green data set.

\begin{table}[t]
  \centering
  \caption{Comparison of methods on the Taylor--Green vortex.  ``Data Error''
    is the mean squared error on the training data after convergence.
    ``Symbolic Output'' indicates whether the method produces a human-readable
    formula.}
  \label{tab:tgv_comparison}
  \smallskip
  \begin{tabular}{@{}lccc@{}}
    \toprule
    Method & Data Error & Symbolic Output & Interpretable \\
    \midrule
    MLP-PINN (3-layer, 64-wide)
      & ${\sim}\,10^{-3}$ & No (weight matrix) & No \\
    Raw EML tree (depth 4)
      & ${\sim}\,0.28$    & Opaque nested \texttt{eml()} & Barely \\
    \textsc{ComposeHead} (free weights)
      & ${\sim}\,10^{-3}$ & Messy multi-term sum & Partially \\
    \textsc{ComposeHead} (separable)
      & $7\!\times\!10^{-8}$ & $\sin(x)\cos(y)\exp(-t/5)$ & \textbf{Yes} \\
    \bottomrule
  \end{tabular}
\end{table}

The MLP-PINN achieves low data error but yields an opaque weight matrix.  The
raw EML tree (depth~4) struggles with the three-variable product structure and
converges to a poor local minimum.  The \textsc{ComposeHead} with free
(non-axis-aligned) weight vectors fits the data reasonably but distributes
signal across multiple terms, producing a multi-term expression that resists
simplification.  Only the separable \textsc{ComposeHead} recovers the exact
symbolic solution.

% ===========================================================================
\subsection{Couette and Poiseuille Flow}
\label{sec:couette_poiseuille}

As elementary validation cases, we apply the method to two one-dimensional
viscous flows for which closed-form solutions are known \emph{a priori}.

\paragraph{Couette flow.}
The steady-state velocity profile between a stationary wall ($y=0$) and a
moving wall ($y=H$, velocity $U_0$) satisfies $\mathrm{d}^2u/\mathrm{d}y^2=0$
with boundary conditions $u(0)=0$, $u(H)=U_0$, giving
$u(y)=U_0\,y/H$.
Using a depth-2 EML tree with PDE-residual loss (enforcing the zero
second derivative) and boundary penalty terms, the method recovers the identity
function $\mathrm{id}(y)$ with the correct multiplicative coefficient.  The
recovery is trivially exact.

\paragraph{Poiseuille flow.}
Pressure-driven flow in a channel of height~$H$ yields
$u(y) = \frac{\Delta P}{2\mu L}(Hy - y^2)$,
a parabolic profile.  Using a depth-3 EML tree with the boundary conditions
$u(0)=u(H)=0$ and the constant-second-derivative constraint
$\mathrm{d}^2u/\mathrm{d}y^2 = -\Delta P / (\mu L)$, the search recovers
$\mathrm{id}(y)$ and $\mathrm{sq}(y)$ terms with the expected coefficients.
The recovery is again trivially exact.

These two cases confirm that the primitive library and snapping pipeline handle
polynomial solutions correctly, complementing the transcendental-function
recovery demonstrated by the Taylor--Green experiment.

% ===========================================================================
\subsection{Algebraic Physical Laws}
\label{sec:algebraic}

To demonstrate that the method generalizes beyond PDE solutions to algebraic
relationships, we apply the raw EML tree search (depth~3--4) to two classical
laws from noisy synthetic data.

\paragraph{Kepler's third law.}
We generate 400 data points $(a_i, T_i)$ with $a\in[0.5, 3.0]$ and
$T = a^{3/2} + \varepsilon$, $\varepsilon\sim\mathcal{N}(0, 0.01)$.  The
search recovers $T = a^{3/2}$ as a composition of EML primitives, with
validation loss below $10^{-4}$ after snapping.

\paragraph{Inverse square law.}
From 400 points $(r_i, F_i)$ with $r\in[0.5, 3.0]$ and
$F = 1/r^{2} + \varepsilon$, $\varepsilon\sim\mathcal{N}(0, 0.01)$, the
search recovers $F = r^{-2}$.

These experiments confirm that EML-based symbolic regression is not restricted
to the separable-product ansatz of the \textsc{ComposeHead} but can also
discover power-law and rational relationships via the general tree search,
provided the target expression is of moderate complexity.

% ===========================================================================
\subsection{Ablation Study}
\label{sec:ablation}

We isolate the effect of each design decision on the Taylor--Green benchmark.
Table~\ref{tab:ablation} summarizes the results.

\begin{table}[t]
  \centering
  \caption{Ablation study on the Taylor--Green vortex.  Each row modifies one
    component of the default pipeline (separable single-term search with
    normalization).  ``Snap Error'' is the mean absolute error of the snapped
    model on 5{,}000 held-out points.}
  \label{tab:ablation}
  \smallskip
  \begin{tabular}{@{}lll@{}}
    \toprule
    Variant & Snap Error & Outcome \\
    \midrule
    \textbf{Full pipeline (default)}
      & $7\!\times\!10^{-8}$ & Exact recovery \\
    Free-weight terms (non-separable)
      & ${\sim}\,10^{-3}$ & Multi-term local minimum \\
    Multi-term with $L_1$ regularization
      & ${\sim}\,10^{-2}$ & Over-pruning of 60+ competing terms \\
    Without normalization (direct snap)
      & $0.005$--$0.02$ & Phase offsets corrupt snapping \\
    \bottomrule
  \end{tabular}
\end{table}

\paragraph{Separable vs.\ free-weight terms.}
When each factor's weight vector is a free $\mathbb{R}^3$ vector rather than
an axis-aligned scalar, the optimizer converges to a representation that
spreads signal across multiple terms.  The resulting multi-term sum achieves
reasonable data error (${\sim}\,10^{-3}$) but does not simplify to the
known single-product expression, because the weight sharing across input
dimensions introduces a degenerate manifold of equivalent parameterizations.

\paragraph{Single-term search vs.\ multi-term with $L_1$.}
A natural alternative to our candidate-by-candidate search is to instantiate
all 27 separable terms simultaneously and apply $L_1$ regularization to the
coefficients to encourage sparsity.  In practice, this approach over-prunes:
with 60+ competing terms (27 for each velocity component plus cross-terms),
the $L_1$ penalty suppresses the correct term before the optimizer can
distinguish it from its neighbors.  Single-term search avoids this failure mode
by giving each candidate the full optimization budget in isolation.

\paragraph{Effect of normalization.}
Without the normalization step, snapping must operate directly on the raw
learned parameters, which carry arbitrary phase offsets
($b \approx -\pi$, $-\pi/2$) and exponential biases ($b \approx 0.6$).
Rounding these values to the nearest ``clean'' constants introduces errors of
$0.005$--$0.02$, two to five orders of magnitude worse than snapping after
normalization.  The normalization step is therefore essential for exact
recovery.

\paragraph{Candidate pool coverage.}
Among the 27 candidates searched for each velocity component, the top~4
candidates all achieve final loss below $10^{-7}$.  These correspond to the
four function triples that are related to the ground truth by the
$\sin\leftrightarrow\cos$ phase-shift equivalence (e.g.,
$\sin\!\cdot\!\sin\!\cdot\!\exp$ and $\cos\!\cdot\!\cos\!\cdot\!\exp$ for the
$u$-component both encode $\sin(x)\cos(y)\exp(-t/5)$ under suitable phase
offsets).  The remaining 23 candidates converge to losses above $10^{-2}$,
providing a clear separation that makes candidate selection unambiguous.
